% =====================================================================
% Empirical Design
% Status: first full draft.
% Each numerical choice that the methodology of
% thesis/methodology.tex defers to this section is recorded
% explicitly here, either as a fixed value or as a precise TODO
% tied to a specific notebook.
% =====================================================================
\section{Empirical Design}\label{sec:empirical-design}

This section translates the methodology of
Section~\ref{sec:methodology} into concrete operational choices for
the KNMI hourly wind gust data. Every numerical choice that the
methodology leaves to the empirical design is either pinned here or
recorded as a TODO that names the notebook responsible for fixing
it.

\subsection{Data paths and input files}\label{sec:emp-paths}

The pipeline reads from two locations under the project root.

\begin{itemize}
    \item Hourly observations:
          \texttt{data/yearly\_aggregated\_FH\_FX/year=YYYY/month=MM/knmi\_fh\_YYYY\_MM.parquet},
          schema
          \texttt{time, station, stationname, lat, lon, height, FH, FX}.
    \item Station metadata:
          \texttt{data/station\_metadata.csv}, schema
          \texttt{station\_id, stationname, lat, lon, height}.
\end{itemize}

The loader is \texttt{python\_project/src/data\_loading.py}. The
station-metadata reader is
\texttt{python\_project/src/station\_metadata.py}.

\subsection{Variable selection}\label{sec:emp-vars}

The only thesis variable is \(\mathrm{FX}\). The hourly mean wind
speed \(\mathrm{FH}\) is read from the same parquet schema but is
not used in any of the dependence estimators. Other columns
(\texttt{stationname}, \texttt{lat}, \texttt{lon}, \texttt{height})
are used for filtering and for the pair distance.

\subsection{Sample period}\label{sec:emp-sample}

\begin{itemize}
    \item Sample frequency: one observation per (station, hour).
    \item Sample window for the main analysis:
          % Data audit 2026-05-26 (outputs/data_audit_report.md):
          %   - 1971 onwards: 13 balanced stations at FX presence
          %     threshold 0.80 across every year.
          %   - 1981 onwards: 14 balanced stations.
          %   - 1991 onwards: 25 balanced stations.
          %   - 2001 onwards: 38 balanced stations.
          %   - 1951 and 1961: only 1 balanced station; FX is
          %     effectively unobserved in the early decades.
          % Recommended candidates: 1991 (good balance of length
          % and panel size) or 1981 (longer window, smaller panel).
          % TODO: pin one start year before estimation.
    \item Sample end year: 2026 (76 calendar years of partitions
          on disk).
          % Note: 2026 is partial -- only 5 months at the audit
          % date. If a complete calendar year is required, use 2025
          % as the end year.
\end{itemize}

\subsection{Station inclusion criteria}\label{sec:emp-stations}

Stations are retained for the main analysis if they satisfy all of
the following.

\begin{enumerate}
    \item Located in the Netherlands (Dutch land or coastal
          stations). North-Sea platform stations are excluded from
          the main panel on the grounds of differing measurement
          environment.
          % Data audit 2026-05-26: 70 unique station ids ever
          % observed across 1951-2026. The balanced panel at
          % start_year=1991, FX-presence threshold 0.80 contains 25
          % stations; the explicit list is in
          % outputs/tables/audit_balanced_panels.csv (column
          % station_ids of the row with
          % start_year=1991, fx_present_share_threshold=0.80).
          % TODO: filter offshore platforms from the 25-station
          % list before adopting it as the main panel.
    \item \(\mathrm{FX}\) sensor documented as representative for
          a reference height of 10\,m.
    \item Present in every winter and every summer of the chosen
          sample window.
    \item Meet the within-season completeness threshold of
          Section~\ref{sec:emp-missing}.
\end{enumerate}

Stations that satisfy criteria 1--2 but fail criterion 3 or 4 are
retained only in the larger \emph{unbalanced panel} used as a
robustness check in Section~\ref{sec:rob-stations}.

\subsection{Season definitions and winter-year
convention}\label{sec:emp-seasons}

\begin{itemize}
    \item Winter \(c = \mathrm{W}\): December, January, February
          (DJF).
    \item Summer \(c = \mathrm{S}\): June, July, August (JJA).
    \item Winter year. December of calendar year \(y\) is assigned
          to winter year \(y+1\). So winter \(1967\) consists of
          December \(1966\), January \(1967\) and February
          \(1967\). This convention is enforced in
          \texttt{python\_project/src/seasons.py}.
\end{itemize}

Spring and autumn hours are dropped.

\subsection{Missingness rules}\label{sec:emp-missing}

\begin{itemize}
    \item Hourly observations with missing \(\mathrm{FX}\) are
          dropped at the hour level.
    \item A block (Section~\ref{sec:meth-extremes}) is discarded
          if more than a fixed fraction of its hours are missing.
          % TODO: set the fraction. Suggested starting value: 0.10.
    \item A station-season-year is discarded if its within-season
          coverage is below a fixed fraction.
          % TODO: set the fraction. Suggested starting value: 0.80.
    \item Station-level pairwise estimators
          (\(\widehat\chi_{ij}^{(c)}\), \(\widehat\theta_{ij}^{(c)}\))
          use pairwise-complete hours.
\end{itemize}

\subsection{Construction of station pairs and
distances}\label{sec:emp-pairs}

The set of unordered station pairs on the balanced panel is
\(\mathcal{P} = \{(i,j) : 1 \le i < j \le N\}\), with cardinality
\(N(N-1)/2\). For each pair the geographical distance is computed
in kilometres by the haversine formula on the WGS84
latitude/longitude coordinates of the two stations:
\[
    d_{ij} \;=\;
    \mathrm{haversine}(s_i, s_j) \quad \text{in km}.
\]
The implementation is in
\texttt{python\_project/src/station\_pairs.py}. As a robustness
check, distances are recomputed in the projected Dutch RD~New
(EPSG:28992) coordinate system; the projected distance differs from
the great-circle distance by at most a few hundred metres for the
Dutch domain and should not affect the dependence-distance summary
materially.

\subsection{Extreme-event construction: numerical
choices}\label{sec:emp-extremes}

\begin{itemize}
    \item Main route (block maxima): block size
          \(n =\)
          % TODO: pin one. Candidates: monthly-within-season (\(n
          %       \approx 720\) hours) or seasonal (\(n \approx
          %       2160\) hours).
          .
    \item Robustness route (POT): threshold quantile
          \(q =\)
          % TODO: pin one. Candidates: 0.95, 0.97, 0.99 of the
          %       within-station-season empirical distribution of
          %       FX.
          .
    \item POT declustering: runs declustering with minimum
          separation
          % TODO: hours. Suggested starting value: 12.
          .
    \item Tail-level grid for
          \(\widehat\chi_{ij}^{(c)}(u)\):
          % TODO: pin a grid. Suggested starting grid:
          %       u in {0.95, 0.96, 0.97, 0.98, 0.99}.
          .
    \item Reference tail level for the headline scalar summary:
          % TODO: pin one u_0, e.g.\ 0.97.
          .
\end{itemize}

\subsection{Distance smoother}\label{sec:emp-smoother}

The dependence--distance smoother in
Section~\ref{sec:meth-distance} is a local-linear kernel smoother
with Gaussian kernel; the bandwidth is selected by leave-one-out
cross-validation per season.
% TODO: confirm bandwidth choice and report it in
%       Section~\ref{sec:results-distance}.

\subsection{Reference distances for scalar
summaries}\label{sec:emp-distances}

The headline scalar summaries
\(\widehat\chi^{(c)}(d_0)\) are reported at a small set of
reference distances \(d_0\). The default set is
\(\{50, 100, 200\}\)\,km.
% TODO: revise the reference set after notebook 01 has produced
%       the empirical distribution of pair distances on the
%       balanced panel.

\subsection{Inference and uncertainty}\label{sec:emp-inference}

\begin{itemize}
    \item Bootstrap method: block bootstrap on the within-season
          hourly panel.
    \item Block length \(L =\)
          % TODO: hours. Should exceed a typical storm duration.
          %       Suggested starting value: 36 hours.
          .
    \item Number of bootstrap replications \(B =\)
          % TODO: pin. Suggested starting value: 1000.
          .
    \item Significance level for confidence bands and for the
          seasonal-equality permutation test: \(\alpha = 0.05\).
\end{itemize}

\subsection{Main sample and robustness samples}\label{sec:emp-samples}

\begin{itemize}
    \item \emph{Main sample.} Balanced station panel
          (\(\S\)~\ref{sec:emp-stations}), block-maxima route
          (\(\S\)~\ref{sec:emp-extremes}), empirical rank-transform
          marginal (\(\S\)~\ref{sec:meth-marginals}), full sample
          window (\(\S\)~\ref{sec:emp-sample}).
    \item \emph{Robustness sample 1.} As main, but unbalanced
          station panel.
    \item \emph{Robustness sample 2.} As main, but threshold-
          exceedance route in place of block maxima.
    \item \emph{Robustness sample 3.} As main, but parametric
          (GEV/GPD) marginal transform.
    \item \emph{Robustness sample 4.} As main, but sample-split
          into an early half and a late half of the sample window.
\end{itemize}

\subsection{Planned outputs and how they enter the
thesis}\label{sec:emp-outputs}

The empirical pipeline is expected to produce the figures and
tables listed below. All outputs are first written to
\texttt{python\_project/output/} and then included in the relevant
thesis section via the standard
\verb|\includegraphics{../python_project/output/figures/...}|
path.

\textbf{Figures} (all to
\texttt{python\_project/output/figures/}):
\begin{enumerate}
    \item \texttt{stations\_map.pdf} — map of the balanced
          station panel. Referenced from
          Section~\ref{sec:data-stations}.
    \item \texttt{station\_coverage\_heatmap.pdf} — station-by-year
          coverage matrix. Referenced from
          Section~\ref{sec:data-descriptives}.
    \item \texttt{seasonal\_maxima\_boxplot.pdf} — per-station
          distributions of seasonal block maxima for winter and
          summer. Referenced from
          Section~\ref{sec:results-marginals}.
    \item \texttt{marginal\_fits\_per\_station.pdf} — GEV/GPD
          diagnostic plots per (station, season).
          Section~\ref{sec:results-marginals}.
    \item \texttt{pair\_chi\_vs\_distance\_by\_season.pdf} —
          pair-level scatter of
          \(\widehat\chi_{ij}^{(c)}\) versus \(d_{ij}\), winter
          and summer panels. Section~\ref{sec:results-pairwise}.
    \item \texttt{chi\_distance\_curves.pdf} — smoothed
          \(\widehat\chi^{(c)}(d)\) and \(\widehat\theta^{(c)}(d)\)
          with bootstrap bands. Section~\ref{sec:results-distance}.
    \item \texttt{chi\_distance\_winter\_minus\_summer.pdf} —
          pointwise difference and bootstrap band, plus
          permutation-test panel.
          Section~\ref{sec:results-comparison}.
\end{enumerate}

\textbf{Tables} (all to \texttt{python\_project/output/tables/}):
\begin{enumerate}
    \item \texttt{station\_inventory.tex} — balanced and unbalanced
          station inventories, with coverage statistics.
          Section~\ref{sec:data-sample}.
    \item \texttt{station\_descriptives.tex} — per-station hourly
          \(\mathrm{FX}\) descriptive statistics, winter vs.\ summer.
          Section~\ref{sec:data-descriptives}.
    \item \texttt{marginal\_fits.tex} — GEV/GPD parameter
          estimates per (station, season).
          Section~\ref{sec:results-marginals}.
    \item \texttt{seasonal\_summary.tex} — scalar dependence
          summaries at the reference distances of
          Section~\ref{sec:emp-distances} with bootstrap intervals.
          Section~\ref{sec:results-comparison}.
\end{enumerate}
